\documentclass{amsart}
\usepackage{amsmath,amssymb,amsthm}

\title{The First-G Law for Semiprimes}
\author{Brayden Ross Sanders / 7Site LLC}
\date{2026}

\newtheorem{theorem}{Theorem}
\newtheorem{corollary}[theorem]{Corollary}
\newtheorem{definition}[theorem]{Definition}

\begin{document}
\maketitle

\begin{abstract}
For a semiprime $b = p \times q$ with $p \leq q$, we prove that the first
non-coprime element in $\{1, \ldots, k\}$ with respect to $b$ appears at
exactly $k = p = \mathrm{SPF}(b)$. Verified across 36,662 $(b,k)$ pairs,
zero exceptions.
\end{abstract}

\section{Definitions}

\begin{definition}[Coprimality partition]
For a semiprime $b$ and $k \geq 1$:
\[
  C(k,b) = \{ x \in \{1,\ldots,k\} : \gcd(x,b)=1 \}, \quad
  G(k,b) = \{ x \in \{1,\ldots,k\} : \gcd(x,b)>1 \}.
\]
\end{definition}

\begin{definition}[First-G event]
$\mathrm{First\text{-}G}(b) = \min\{k \geq 1 : |G(k,b)| > 0\}$.
\end{definition}

\section{Main Theorem}

\begin{theorem}[First-G Law]
For every semiprime $b = p \times q$ with $p \leq q$:
\[
  |G(k,b)| = 0 \quad \text{for all } k < p, \qquad
  |G(p,b)| = 1 \quad (\text{exactly the element } p).
\]
Equivalently, $\mathrm{First\text{-}G}(b) = p = \mathrm{SPF}(b)$.
\end{theorem}

\begin{proof}
Let $x \in \{1,\ldots,k\}$ with $k < p$.
Since $p$ and $q$ are the only prime factors of $b = p \cdot q$,
and since $x < p \leq q$: $x$ is not divisible by $p$ (as $p$ is prime and $x < p$),
and $x$ is not divisible by $q$ (as $q \geq p > x$).
Therefore $\gcd(x, p \cdot q) = 1$, so $x \in C(k,b)$.
Since $x$ was arbitrary, $|G(k,b)| = 0$ for all $k < p$.

At $k = p$: the element $p$ enters $\{1,\ldots,k\}$.
Since $p \mid b$, we have $\gcd(p,b) = p > 1$, so $p \in G(p,b)$.
No element of $\{1,\ldots,p-1\}$ is in $G$ (just proved).
Therefore $|G(p,b)| = 1$.
\end{proof}

\begin{corollary}[Forbidden zone]
The set $\{1, \ldots, p-1\}$ is a \emph{forbidden zone} of width $p-1$:
no non-trivial structure with respect to $b$ exists below position $k = p$.
\end{corollary}

\begin{corollary}[Perfect square semiprimes]
The same argument holds for $b = p^2$ (taking $p = q$).
\end{corollary}

\end{document}
